Most \lora recipes hard-code a rank of 8 or 16. A common belief holds that this default suits easy tasks such as sentiment classification but under-serves generative tasks such as math and code, which are thought to need rank 32 or more. We test this belief with a controlled sweep of about 200 adapters: six tasks (four classification tasks, \gsm, and Python code) by eight ranks (1--128) on \qsmall (two seeds) and \qlarge, with a fixed data budget of 4,000 examples, plus a learning-rate sweep and a cheap predictor based on the spectrum of the full-weight gradient at initialization. At a fixed learning rate of $2\times10^{-4}$, the best rank does vary from 1 to 128, but in the opposite direction from the belief. \gsm and code do best at low rank and degrade at rank 128 (\gsm: 37.3\% at $r{=}1$ vs.\ 30.1\% at $r{=}128$). The only task where rank 8 loses significantly after Holm correction is 6-way emotion classification ($+2.6$\pp at 0.5B, $+2.2$\pp at 1.5B), and this gap disappears once the learning rate is tuned per rank: rank 8 reaches 92.8\%, matching ranks 32 and 128. The gradient-spectrum statistic orders tasks consistently across model sizes, but it is \emph{negatively} correlated with the empirical knee rank (Spearman $\rho=-0.74$, $p=0.014$, $n=12$). In the small-data supervised fine-tuning regime we study, rank 8 is not a blind spot. Pairing a large rank with an untuned learning rate is the real risk.
